\documentclass{article}
\usepackage{icomp2026_conference,times}
\usepackage{amsmath}
\usepackage{amssymb}
\usepackage{amsfonts}
\usepackage{graphicx}
\usepackage{booktabs}
\usepackage{bbold}
\usepackage{amsthm}
\usepackage{float}
%%%%% NEW MATH DEFINITIONS %%%%%

\usepackage{amsmath,amsfonts,bm}

% Mark sections of captions for referring to divisions of figures
\newcommand{\figleft}{{\em (Left)}}
\newcommand{\figcenter}{{\em (Center)}}
\newcommand{\figright}{{\em (Right)}}
\newcommand{\figtop}{{\em (Top)}}
\newcommand{\figbottom}{{\em (Bottom)}}
\newcommand{\captiona}{{\em (a)}}
\newcommand{\captionb}{{\em (b)}}
\newcommand{\captionc}{{\em (c)}}
\newcommand{\captiond}{{\em (d)}}

% Highlight a newly defined term
\newcommand{\newterm}[1]{{\bf #1}}


% Figure reference, lower-case.
\def\figref#1{figure~\ref{#1}}
% Figure reference, capital. For start of sentence
\def\Figref#1{Figure~\ref{#1}}
\def\twofigref#1#2{figures \ref{#1} and \ref{#2}}
\def\quadfigref#1#2#3#4{figures \ref{#1}, \ref{#2}, \ref{#3} and \ref{#4}}
% Section reference, lower-case.
\def\secref#1{section~\ref{#1}}
% Section reference, capital.
\def\Secref#1{Section~\ref{#1}}
% Reference to two sections.
\def\twosecrefs#1#2{sections \ref{#1} and \ref{#2}}
% Reference to three sections.
\def\secrefs#1#2#3{sections \ref{#1}, \ref{#2} and \ref{#3}}
% Reference to an equation, lower-case.
\def\eqref#1{equation~\ref{#1}}
% Reference to an equation, upper case
\def\Eqref#1{Equation~\ref{#1}}
% A raw reference to an equation---avoid using if possible
\def\plaineqref#1{\ref{#1}}
% Reference to a chapter, lower-case.
\def\chapref#1{chapter~\ref{#1}}
% Reference to an equation, upper case.
\def\Chapref#1{Chapter~\ref{#1}}
% Reference to a range of chapters
\def\rangechapref#1#2{chapters\ref{#1}--\ref{#2}}
% Reference to an algorithm, lower-case.
\def\algref#1{algorithm~\ref{#1}}
% Reference to an algorithm, upper case.
\def\Algref#1{Algorithm~\ref{#1}}
\def\twoalgref#1#2{algorithms \ref{#1} and \ref{#2}}
\def\Twoalgref#1#2{Algorithms \ref{#1} and \ref{#2}}
% Reference to a part, lower case
\def\partref#1{part~\ref{#1}}
% Reference to a part, upper case
\def\Partref#1{Part~\ref{#1}}
\def\twopartref#1#2{parts \ref{#1} and \ref{#2}}

\def\ceil#1{\lceil #1 \rceil}
\def\floor#1{\lfloor #1 \rfloor}
\def\1{\bm{1}}
\newcommand{\train}{\mathcal{D}}
\newcommand{\valid}{\mathcal{D_{\mathrm{valid}}}}
\newcommand{\test}{\mathcal{D_{\mathrm{test}}}}

\def\eps{{\epsilon}}


% Random variables
\def\reta{{\textnormal{$\eta$}}}
\def\ra{{\textnormal{a}}}
\def\rb{{\textnormal{b}}}
\def\rc{{\textnormal{c}}}
\def\rd{{\textnormal{d}}}
\def\re{{\textnormal{e}}}
\def\rf{{\textnormal{f}}}
\def\rg{{\textnormal{g}}}
\def\rh{{\textnormal{h}}}
\def\ri{{\textnormal{i}}}
\def\rj{{\textnormal{j}}}
\def\rk{{\textnormal{k}}}
\def\rl{{\textnormal{l}}}
% rm is already a command, just don't name any random variables m
\def\rn{{\textnormal{n}}}
\def\ro{{\textnormal{o}}}
\def\rp{{\textnormal{p}}}
\def\rq{{\textnormal{q}}}
\def\rr{{\textnormal{r}}}
\def\rs{{\textnormal{s}}}
\def\rt{{\textnormal{t}}}
\def\ru{{\textnormal{u}}}
\def\rv{{\textnormal{v}}}
\def\rw{{\textnormal{w}}}
\def\rx{{\textnormal{x}}}
\def\ry{{\textnormal{y}}}
\def\rz{{\textnormal{z}}}

% Random vectors
\def\rvepsilon{{\mathbf{\epsilon}}}
\def\rvtheta{{\mathbf{\theta}}}
\def\rva{{\mathbf{a}}}
\def\rvb{{\mathbf{b}}}
\def\rvc{{\mathbf{c}}}
\def\rvd{{\mathbf{d}}}
\def\rve{{\mathbf{e}}}
\def\rvf{{\mathbf{f}}}
\def\rvg{{\mathbf{g}}}
\def\rvh{{\mathbf{h}}}
\def\rvu{{\mathbf{i}}}
\def\rvj{{\mathbf{j}}}
\def\rvk{{\mathbf{k}}}
\def\rvl{{\mathbf{l}}}
\def\rvm{{\mathbf{m}}}
\def\rvn{{\mathbf{n}}}
\def\rvo{{\mathbf{o}}}
\def\rvp{{\mathbf{p}}}
\def\rvq{{\mathbf{q}}}
\def\rvr{{\mathbf{r}}}
\def\rvs{{\mathbf{s}}}
\def\rvt{{\mathbf{t}}}
\def\rvu{{\mathbf{u}}}
\def\rvv{{\mathbf{v}}}
\def\rvw{{\mathbf{w}}}
\def\rvx{{\mathbf{x}}}
\def\rvy{{\mathbf{y}}}
\def\rvz{{\mathbf{z}}}

% Elements of random vectors
\def\erva{{\textnormal{a}}}
\def\ervb{{\textnormal{b}}}
\def\ervc{{\textnormal{c}}}
\def\ervd{{\textnormal{d}}}
\def\erve{{\textnormal{e}}}
\def\ervf{{\textnormal{f}}}
\def\ervg{{\textnormal{g}}}
\def\ervh{{\textnormal{h}}}
\def\ervi{{\textnormal{i}}}
\def\ervj{{\textnormal{j}}}
\def\ervk{{\textnormal{k}}}
\def\ervl{{\textnormal{l}}}
\def\ervm{{\textnormal{m}}}
\def\ervn{{\textnormal{n}}}
\def\ervo{{\textnormal{o}}}
\def\ervp{{\textnormal{p}}}
\def\ervq{{\textnormal{q}}}
\def\ervr{{\textnormal{r}}}
\def\ervs{{\textnormal{s}}}
\def\ervt{{\textnormal{t}}}
\def\ervu{{\textnormal{u}}}
\def\ervv{{\textnormal{v}}}
\def\ervw{{\textnormal{w}}}
\def\ervx{{\textnormal{x}}}
\def\ervy{{\textnormal{y}}}
\def\ervz{{\textnormal{z}}}

% Random matrices
\def\rmA{{\mathbf{A}}}
\def\rmB{{\mathbf{B}}}
\def\rmC{{\mathbf{C}}}
\def\rmD{{\mathbf{D}}}
\def\rmE{{\mathbf{E}}}
\def\rmF{{\mathbf{F}}}
\def\rmG{{\mathbf{G}}}
\def\rmH{{\mathbf{H}}}
\def\rmI{{\mathbf{I}}}
\def\rmJ{{\mathbf{J}}}
\def\rmK{{\mathbf{K}}}
\def\rmL{{\mathbf{L}}}
\def\rmM{{\mathbf{M}}}
\def\rmN{{\mathbf{N}}}
\def\rmO{{\mathbf{O}}}
\def\rmP{{\mathbf{P}}}
\def\rmQ{{\mathbf{Q}}}
\def\rmR{{\mathbf{R}}}
\def\rmS{{\mathbf{S}}}
\def\rmT{{\mathbf{T}}}
\def\rmU{{\mathbf{U}}}
\def\rmV{{\mathbf{V}}}
\def\rmW{{\mathbf{W}}}
\def\rmX{{\mathbf{X}}}
\def\rmY{{\mathbf{Y}}}
\def\rmZ{{\mathbf{Z}}}

% Elements of random matrices
\def\ermA{{\textnormal{A}}}
\def\ermB{{\textnormal{B}}}
\def\ermC{{\textnormal{C}}}
\def\ermD{{\textnormal{D}}}
\def\ermE{{\textnormal{E}}}
\def\ermF{{\textnormal{F}}}
\def\ermG{{\textnormal{G}}}
\def\ermH{{\textnormal{H}}}
\def\ermI{{\textnormal{I}}}
\def\ermJ{{\textnormal{J}}}
\def\ermK{{\textnormal{K}}}
\def\ermL{{\textnormal{L}}}
\def\ermM{{\textnormal{M}}}
\def\ermN{{\textnormal{N}}}
\def\ermO{{\textnormal{O}}}
\def\ermP{{\textnormal{P}}}
\def\ermQ{{\textnormal{Q}}}
\def\ermR{{\textnormal{R}}}
\def\ermS{{\textnormal{S}}}
\def\ermT{{\textnormal{T}}}
\def\ermU{{\textnormal{U}}}
\def\ermV{{\textnormal{V}}}
\def\ermW{{\textnormal{W}}}
\def\ermX{{\textnormal{X}}}
\def\ermY{{\textnormal{Y}}}
\def\ermZ{{\textnormal{Z}}}

% Vectors
\def\vzero{{\bm{0}}}
\def\vone{{\bm{1}}}
\def\vmu{{\bm{\mu}}}
\def\vtheta{{\bm{\theta}}}
\def\va{{\bm{a}}}
\def\vb{{\bm{b}}}
\def\vc{{\bm{c}}}
\def\vd{{\bm{d}}}
\def\ve{{\bm{e}}}
\def\vf{{\bm{f}}}
\def\vg{{\bm{g}}}
\def\vh{{\bm{h}}}
\def\vi{{\bm{i}}}
\def\vj{{\bm{j}}}
\def\vk{{\bm{k}}}
\def\vl{{\bm{l}}}
\def\vm{{\bm{m}}}
\def\vn{{\bm{n}}}
\def\vo{{\bm{o}}}
\def\vp{{\bm{p}}}
\def\vq{{\bm{q}}}
\def\vr{{\bm{r}}}
\def\vs{{\bm{s}}}
\def\vt{{\bm{t}}}
\def\vu{{\bm{u}}}
\def\vv{{\bm{v}}}
\def\vw{{\bm{w}}}
\def\vx{{\bm{x}}}
\def\vy{{\bm{y}}}
\def\vz{{\bm{z}}}

% Elements of vectors
\def\evalpha{{\alpha}}
\def\evbeta{{\beta}}
\def\evepsilon{{\epsilon}}
\def\evlambda{{\lambda}}
\def\evomega{{\omega}}
\def\evmu{{\mu}}
\def\evpsi{{\psi}}
\def\evsigma{{\sigma}}
\def\evtheta{{\theta}}
\def\eva{{a}}
\def\evb{{b}}
\def\evc{{c}}
\def\evd{{d}}
\def\eve{{e}}
\def\evf{{f}}
\def\evg{{g}}
\def\evh{{h}}
\def\evi{{i}}
\def\evj{{j}}
\def\evk{{k}}
\def\evl{{l}}
\def\evm{{m}}
\def\evn{{n}}
\def\evo{{o}}
\def\evp{{p}}
\def\evq{{q}}
\def\evr{{r}}
\def\evs{{s}}
\def\evt{{t}}
\def\evu{{u}}
\def\evv{{v}}
\def\evw{{w}}
\def\evx{{x}}
\def\evy{{y}}
\def\evz{{z}}

% Matrix
\def\mA{{\bm{A}}}
\def\mB{{\bm{B}}}
\def\mC{{\bm{C}}}
\def\mD{{\bm{D}}}
\def\mE{{\bm{E}}}
\def\mF{{\bm{F}}}
\def\mG{{\bm{G}}}
\def\mH{{\bm{H}}}
\def\mI{{\bm{I}}}
\def\mJ{{\bm{J}}}
\def\mK{{\bm{K}}}
\def\mL{{\bm{L}}}
\def\mM{{\bm{M}}}
\def\mN{{\bm{N}}}
\def\mO{{\bm{O}}}
\def\mP{{\bm{P}}}
\def\mQ{{\bm{Q}}}
\def\mR{{\bm{R}}}
\def\mS{{\bm{S}}}
\def\mT{{\bm{T}}}
\def\mU{{\bm{U}}}
\def\mV{{\bm{V}}}
\def\mW{{\bm{W}}}
\def\mX{{\bm{X}}}
\def\mY{{\bm{Y}}}
\def\mZ{{\bm{Z}}}
\def\mBeta{{\bm{\beta}}}
\def\mPhi{{\bm{\Phi}}}
\def\mLambda{{\bm{\Lambda}}}
\def\mSigma{{\bm{\Sigma}}}

% Tensor
\DeclareMathAlphabet{\mathsfit}{\encodingdefault}{\sfdefault}{m}{sl}
\SetMathAlphabet{\mathsfit}{bold}{\encodingdefault}{\sfdefault}{bx}{n}
\newcommand{\tens}[1]{\bm{\mathsfit{#1}}}
\def\tA{{\tens{A}}}
\def\tB{{\tens{B}}}
\def\tC{{\tens{C}}}
\def\tD{{\tens{D}}}
\def\tE{{\tens{E}}}
\def\tF{{\tens{F}}}
\def\tG{{\tens{G}}}
\def\tH{{\tens{H}}}
\def\tI{{\tens{I}}}
\def\tJ{{\tens{J}}}
\def\tK{{\tens{K}}}
\def\tL{{\tens{L}}}
\def\tM{{\tens{M}}}
\def\tN{{\tens{N}}}
\def\tO{{\tens{O}}}
\def\tP{{\tens{P}}}
\def\tQ{{\tens{Q}}}
\def\tR{{\tens{R}}}
\def\tS{{\tens{S}}}
\def\tT{{\tens{T}}}
\def\tU{{\tens{U}}}
\def\tV{{\tens{V}}}
\def\tW{{\tens{W}}}
\def\tX{{\tens{X}}}
\def\tY{{\tens{Y}}}
\def\tZ{{\tens{Z}}}


% Graph
\def\gA{{\mathcal{A}}}
\def\gB{{\mathcal{B}}}
\def\gC{{\mathcal{C}}}
\def\gD{{\mathcal{D}}}
\def\gE{{\mathcal{E}}}
\def\gF{{\mathcal{F}}}
\def\gG{{\mathcal{G}}}
\def\gH{{\mathcal{H}}}
\def\gI{{\mathcal{I}}}
\def\gJ{{\mathcal{J}}}
\def\gK{{\mathcal{K}}}
\def\gL{{\mathcal{L}}}
\def\gM{{\mathcal{M}}}
\def\gN{{\mathcal{N}}}
\def\gO{{\mathcal{O}}}
\def\gP{{\mathcal{P}}}
\def\gQ{{\mathcal{Q}}}
\def\gR{{\mathcal{R}}}
\def\gS{{\mathcal{S}}}
\def\gT{{\mathcal{T}}}
\def\gU{{\mathcal{U}}}
\def\gV{{\mathcal{V}}}
\def\gW{{\mathcal{W}}}
\def\gX{{\mathcal{X}}}
\def\gY{{\mathcal{Y}}}
\def\gZ{{\mathcal{Z}}}

% Sets
\def\sA{{\mathbb{A}}}
\def\sB{{\mathbb{B}}}
\def\sC{{\mathbb{C}}}
\def\sD{{\mathbb{D}}}
% Don't use a set called E, because this would be the same as our symbol
% for expectation.
\def\sF{{\mathbb{F}}}
\def\sG{{\mathbb{G}}}
\def\sH{{\mathbb{H}}}
\def\sI{{\mathbb{I}}}
\def\sJ{{\mathbb{J}}}
\def\sK{{\mathbb{K}}}
\def\sL{{\mathbb{L}}}
\def\sM{{\mathbb{M}}}
\def\sN{{\mathbb{N}}}
\def\sO{{\mathbb{O}}}
\def\sP{{\mathbb{P}}}
\def\sQ{{\mathbb{Q}}}
\def\sR{{\mathbb{R}}}
\def\sS{{\mathbb{S}}}
\def\sT{{\mathbb{T}}}
\def\sU{{\mathbb{U}}}
\def\sV{{\mathbb{V}}}
\def\sW{{\mathbb{W}}}
\def\sX{{\mathbb{X}}}
\def\sY{{\mathbb{Y}}}
\def\sZ{{\mathbb{Z}}}

% Entries of a matrix
\def\emLambda{{\Lambda}}
\def\emA{{A}}
\def\emB{{B}}
\def\emC{{C}}
\def\emD{{D}}
\def\emE{{E}}
\def\emF{{F}}
\def\emG{{G}}
\def\emH{{H}}
\def\emI{{I}}
\def\emJ{{J}}
\def\emK{{K}}
\def\emL{{L}}
\def\emM{{M}}
\def\emN{{N}}
\def\emO{{O}}
\def\emP{{P}}
\def\emQ{{Q}}
\def\emR{{R}}
\def\emS{{S}}
\def\emT{{T}}
\def\emU{{U}}
\def\emV{{V}}
\def\emW{{W}}
\def\emX{{X}}
\def\emY{{Y}}
\def\emZ{{Z}}
\def\emSigma{{\Sigma}}

% entries of a tensor
% Same font as tensor, without \bm wrapper
\newcommand{\etens}[1]{\mathsfit{#1}}
\def\etLambda{{\etens{\Lambda}}}
\def\etA{{\etens{A}}}
\def\etB{{\etens{B}}}
\def\etC{{\etens{C}}}
\def\etD{{\etens{D}}}
\def\etE{{\etens{E}}}
\def\etF{{\etens{F}}}
\def\etG{{\etens{G}}}
\def\etH{{\etens{H}}}
\def\etI{{\etens{I}}}
\def\etJ{{\etens{J}}}
\def\etK{{\etens{K}}}
\def\etL{{\etens{L}}}
\def\etM{{\etens{M}}}
\def\etN{{\etens{N}}}
\def\etO{{\etens{O}}}
\def\etP{{\etens{P}}}
\def\etQ{{\etens{Q}}}
\def\etR{{\etens{R}}}
\def\etS{{\etens{S}}}
\def\etT{{\etens{T}}}
\def\etU{{\etens{U}}}
\def\etV{{\etens{V}}}
\def\etW{{\etens{W}}}
\def\etX{{\etens{X}}}
\def\etY{{\etens{Y}}}
\def\etZ{{\etens{Z}}}

% The true underlying data generating distribution
\newcommand{\pdata}{p_{\rm{data}}}
% The empirical distribution defined by the training set
\newcommand{\ptrain}{\hat{p}_{\rm{data}}}
\newcommand{\Ptrain}{\hat{P}_{\rm{data}}}
% The model distribution
\newcommand{\pmodel}{p_{\rm{model}}}
\newcommand{\Pmodel}{P_{\rm{model}}}
\newcommand{\ptildemodel}{\tilde{p}_{\rm{model}}}
% Stochastic autoencoder distributions
\newcommand{\pencode}{p_{\rm{encoder}}}
\newcommand{\pdecode}{p_{\rm{decoder}}}
\newcommand{\precons}{p_{\rm{reconstruct}}}

\newcommand{\laplace}{\mathrm{Laplace}} % Laplace distribution

\newcommand{\E}{\mathbb{E}}
\newcommand{\Ls}{\mathcal{L}}
\newcommand{\R}{\mathbb{R}}
\newcommand{\emp}{\tilde{p}}
\newcommand{\lr}{\alpha}
\newcommand{\reg}{\lambda}
\newcommand{\rect}{\mathrm{rectifier}}
\newcommand{\softmax}{\mathrm{softmax}}
\newcommand{\sigmoid}{\sigma}
\newcommand{\softplus}{\zeta}
\newcommand{\KL}{D_{\mathrm{KL}}}
\newcommand{\Var}{\mathrm{Var}}
\newcommand{\standarderror}{\mathrm{SE}}
\newcommand{\Cov}{\mathrm{Cov}}
% Wolfram Mathworld says $L^2$ is for function spaces and $\ell^2$ is for vectors
% But then they seem to use $L^2$ for vectors throughout the site, and so does
% wikipedia.
\newcommand{\normlzero}{L^0}
\newcommand{\normlone}{L^1}
\newcommand{\normltwo}{L^2}
\newcommand{\normlp}{L^p}
\newcommand{\normmax}{L^\infty}

\newcommand{\parents}{Pa} % See usage in notation.tex. Chosen to match Daphne's book.

\DeclareMathOperator*{\argmax}{arg\,max}
\DeclareMathOperator*{\argmin}{arg\,min}

\DeclareMathOperator{\sign}{sign}
\DeclareMathOperator{\Tr}{Tr}
\let\ab\allowbreak


\usepackage{hyperref}
\hypersetup{hidelinks}
\usepackage{url}
\usepackage{cleveref}

\title{Min-Marginal Power Indices for \\ Multi-Objective Feature Selection}

\author{Anonymous Authors}

\begin{document}
\maketitle

\begin{abstract}
Game-theoretic feature selection methods score features by their mean marginal contribution across coalitions, concealing a risk: a feature beneficial on average may be harmful in specific contexts. We introduce Min-Marginal Power Indices (MinPPI), which replace the mean with the minimum under a binary utility-threshold voting game $v_\alpha(S) = \mathbb{1}[U(\mathcal{M}_S) \geq \alpha \cdot U_{\text{full}}]$ for a configured higher-is-better utility $U$. This min--binary pairing---an unoccupied cell in the design space of game-theoretic feature selectors---yields a discrete $\{-1, 0, +1\}$ output that classifies each feature as harmful, neutral, or decisive with no arbitrary cutoff. The two power-index formulations (Min Shapley-Shubik and Min Banzhaf) coincide. We develop an MI-correlation-guided adversarial search (Algorithm~B) and give a conditional geometric detection bound for harmful witnesses in the perturbation distribution (Theorem~2). MinPPI constructs $\alpha$-parameterized (compression, utility preservation) operating points via parametric scalarization, requiring no evolutionary search.

The empirical component is a pre-registered protocol rather than a completed study. It covers seven tabular datasets, five canonical seeds, matched removal controls, an $\alpha$ sweep, estimator and utility ablations, and a full-versus-prefiltered NSGA-II comparison. No formal outcome is reported in this draft: empirical claims will be added only from a newly completed provenance bundle whose resolved protocol, complete upstream artifacts, canonical summary tables, and publication eligibility have been validated.
\end{abstract}

\section{Introduction}
\label{sec:intro}

Feature selection identifies compact yet predictive feature subsets and is inherently multi-objective: practitioners balance cardinality, accuracy, and training cost~\citep{guyon2003introduction,chandrashekar2014survey}. Cooperative game theory offers an axiomatically grounded approach: features are players whose coalition value is model performance~\citep{cohen2007feature}. The Shapley value~\citep{shapley1953value} and Banzhaf index~\citep{banzhaf1965weighted} quantify mean marginal contributions; SHAP~\citep{lundberg2017unified} brought these to mainstream ML. Power-index feature selectors include Banzhaf-FS~\citep{karavas2017feature}, Powershap~\citep{verhaeghe2022powershap}, and ShapG~\citep{zhao2025shapg}. Multi-objective evolutionary methods, notably NSGA-II~\citep{deb2002fast,hamdani2007multi}, recover full Pareto fronts but treat features symmetrically during search.

The mean operator has a blind spot: a feature beneficial on average may be actively harmful in specific coalitions. MinShap~\citep{zheng2026minshap} replaces the mean with the minimum, exposing detrimental features, but uses a continuous value function without a principled decision boundary---users must choose arbitrary cutoffs on real-valued scores. We address this gap by pairing the minimum with a \emph{binary} utility-threshold formulation~\citep{owen1972multilinear}: $v_\alpha(S) = \mathbb{1}[U(\mathcal{M}_S) \geq \alpha \cdot U_{\text{full}}]$ for any higher-is-better utility $U$. The min--binary pairing yields a discrete output in $\{-1, 0, +1\}$ and a \emph{threshold-free} three-way classification: features are harmful (there exists a coalition where adding them flips utility below $\alpha$), neutral, or decisive (every coalition requires them). No existing method occupies this cell of the design space (min operator $\times$ binary value function; Table~\ref{tab:op-table} in Appendix).

Our contributions are: (1)~Min-Marginal Power Indices (MinPPI) on the utility-threshold voting game (UTVG), with the property that MinSS and MinBz coincide and give threshold-free three-way classification; (2)~MI-correlation-guided adversarial search (Algorithm~B) with a conditional harmful-witness detection bound (Theorem~2); (3)~an $\alpha$-parameterized multi-objective framework that sweeps the quota $\alpha$ to construct (compression, utility preservation) operating points without evolutionary search; and (4)~a pre-registered empirical protocol on seven tabular datasets that evaluates matched-budget removal, estimator behavior, quota sensitivity, and full-versus-prefiltered NSGA-II. The fourth contribution is a study design; its formal results await a new validated provenance bundle.

\section{Method}
\label{sec:method}

\subsection{Utility-Threshold Voting Game (UTVG)}

Let $\mathcal{D}$ be a tabular dataset with feature set $N = \{f_1, \dots, f_n\}$. For $S \subseteq N$, $\mathcal{M}_S$ is a base learner trained on features $S$, and $U(\mathcal{M}_S)$ is a configured higher-is-better utility. The formal protocol uses ROC AUC for binary classification, one-vs-rest (OVR) AUC for ISOLET multiclass classification, and $R^2$ for regression. The full-feature utility $U_{\text{full}} = U(\mathcal{M}_N)$ must be positive. For quota $\alpha \in (0,1]$, the UTVG is:
\begin{equation}
v_\alpha(S) = \mathbb{1}\!\bigl[\,U(\mathcal{M}_S) \;\geq\; \alpha \cdot U_{\text{full}}\,\bigr].
\label{eq:valpha}
\end{equation}
$v_\alpha$ is evaluated by training $\mathcal{M}_S$---there is no additive surrogate. Unlike classical weighted voting games, $v_\alpha$ is non-monotone: adding features can decrease utility (noise, redundancy), so $S \subset T$ with $v_\alpha(S)=1, v_\alpha(T)=0$ is possible. This non-monotonicity is the empirical basis for detecting harmful features.

\subsection{Min-Marginal Power Indices}

\textbf{Definition 1 (Min Shapley-Shubik, MinSS).}
$\text{MinSS}_i = \min_{\pi \in \Pi(N)} [\, v_\alpha(S_\pi^i \cup \{i\}) - v_\alpha(S_\pi^i) \,]$, where $S_\pi^i$ are features preceding $i$ in permutation $\pi$.

\textbf{Definition 2 (Min Banzhaf, MinBz).}
$\text{MinBz}_i = \min_{S \subseteq N \setminus \{i\}} [v_\alpha(S \cup \{i\}) - v_\alpha(S)]$.

\textbf{Property 1 (Trinary Range).} $\text{MinSS}_i, \text{MinBz}_i \in \{-1, 0, +1\}$ (immediate from $v_\alpha \in \{0,1\}$).

\textbf{Property 2 (Three-Way Classification).} Under $(N, v_\alpha)$: (a)~$\text{MinSS}_i = +1$ (decisive) iff $\forall S: v_\alpha(S)=0, v_\alpha(S\cup\{i\})=1$; (b)~$\text{MinSS}_i = -1$ (harmful) iff $\exists S: v_\alpha(S)=1, v_\alpha(S\cup\{i\})=0$; (c)~$\text{MinSS}_i = 0$ (neutral) otherwise. A feature ``in every winning coalition'' need not be decisive---it can be neutral when its removal never creates a new winning coalition.

\textbf{Property 3 (MinSS = MinBz).} Both minima range over the same distinct marginals: every $S \subseteq N \setminus \{i\}$ occurs as the predecessor set of $i$ in at least one permutation, and repeated occurrences do not affect a minimum. We unify them as $\text{MinPPI}_i$.

The sign identifies $-1$ features as candidates for removal and $0$ features for downstream analysis; actual removal is subject to the pre-registered validation non-inferiority gate in Section~\ref{sec:protocol}. The three-way classification is threshold-free conditional on the configured $\alpha$: boundaries at $0$ and $\pm 1$ follow from $v_\alpha \in \{0,1\}$ with no further discretization step. Proofs are deferred to Appendix~\ref{app:proofs}.

\subsection{Estimation Algorithms}

Exact MinPPI requires $O(n\cdot 2^{n-1})$ evaluations. We give three approximation strategies.

\subsubsection{Algorithm~A (Random MC)}

Draw $K$ random subsets (size-uniform stratified). Return $\widehat{\text{MinPPI}}_i = \min_k [v_\alpha(S^{(k)}\cup\{i\}) - v_\alpha(S^{(k)})]$. Lemma~1 (Appendix~\ref{app:proofs}): $\widehat{\text{MinPPI}}_i^{\text{MC}} \geq \text{MinPPI}_i$ almost surely (conservative: never falsely flags a neutral feature as harmful). If $\text{MinPPI}_i = -1$ and $p_-^i$ is the probability of a harmful marginal, $\Pr[\text{miss}] = (1-p_-^i)^K$. For sparse harmful events ($p_-^i \ll 1$), $K = \Omega(1/p_-)$ is required.

\subsubsection{Algorithm~B (MI-Correlation-Guided Adversarial Search)}

When harmful marginals are sparse, random sampling is inefficient. Harmful coalitions often concentrate where feature $i$ conflicts with correlated neighbors. Algorithm~B: (1)~build a maximum-weight spanning tree $G_{\text{corr}}$ of absolute Pearson correlations; (2)~rank features by mutual information (classification or regression); (3)~for each feature $i$, seed $S_{\text{seed}}$ from its strongest incident neighbors in $G_{\text{corr}}$ plus top mutual-information features; (4)~perturb $S_{\text{seed}}$ for $K_{\text{adv}}$ rounds, returning the minimum observed marginal.

\textbf{Theorem 2 (Conditional Detection).} Let $P_i$ be the coalition distribution induced for feature $i$ by Algorithm~B's complete perturbation step, including both deletion from the seed and addition of features outside it. If a harmful witness $S_i^*$ has probability mass $\rho_i=P_i(S_i^*)>0$, and the $K_{\text{adv}}$ perturbations are independent draws from $P_i$, then $\Pr[\text{detection}] \geq 1-(1-\rho_i)^{K_{\text{adv}}}$. Positive mass is an explicit support assumption; merely having $S_i^*\subseteq S_{\text{seed}}$ is not sufficient because the addition step may also change the coalition. Ground-truth witnesses and their probability mass are unknown on real data, so the theorem does not assert a dataset-independent sample complexity and Algorithm~B remains a heuristic whose advantage must be measured by the formal guidance ablation.

\subsubsection{Estimator Dispatch}

The formal estimator stage compares exact enumeration, random MC, MI-correlation-guided search, and the hybrid dispatcher at the dataset-specific budgets frozen in the resolved protocol. Exact enumeration is permitted only up to 16 selected features and therefore applies to Abalone in the registered datasets. The hybrid combines the configured approximate estimators without changing their frozen coalition requests. The correlation graph and spanning tree are built once at $O(n^2\log n)$ cost.

\subsection{$\alpha$-Parameterized Multi-Objective Selection}
\label{sec:alpha-pareto}

MinPPI constructs multi-objective trade-off points via parametric scalarization---the same principle as the $\varepsilon$-constraint method~\citep{haimes1971bicriterion}. For a fixed dataset and quota $\alpha$, estimation and retraining yield one (compression rate, utility preservation) operating point. Varying $\alpha \in \{0.85,0.90,0.95,0.99\}$ changes which fitted coalitions meet the threshold. Because the learned utility is non-monotone and estimation is finite-budget, the protocol does not impose a monotonic compression or preservation trend across this grid. Every estimated harmful label is accompanied by an observed coalition where $v_\alpha(S)=1$ but $v_\alpha(S\cup\{i\})=0$; whether that label is retained for removal is decided by the validation gate.

\section{Experiments}
\label{sec:experiments}

\subsection{Setup}

The empirical study is pre-registered over seven tabular datasets after training-only variance filtering and mutual-information selection to at most 100 features (Table~\ref{tab:datasets}). Each seed uses a 64/16/20 train/validation/test split, stratified for classification. Preprocessing choices, including feature selection and standardization, are fitted on the training partition only. The canonical seeds are 42, 123, 456, 789, and 1024.

Every task uses XGBoost with a configured reference and cap of 100 boosting rounds, maximum depth 4, learning rate 0.1, and \texttt{n\_jobs=1}. Before any formal test evaluation, training/validation-only calibration selects one dataset-specific round count from $\{10,20,40,100\}$; the resolved model policy freezes that choice, and every downstream variant for the dataset reuses it. Classification models produce probabilities. The dataset configuration selects the higher-is-better utility: ROC AUC for each binary dataset, OVR AUC for ISOLET, and $R^2$ for both regression datasets. Accuracy appears only in the ISOLET utility ablation and is not the formal ISOLET utility. The portable configuration defaults to the CPU histogram backend. A formal campaign may explicitly override the backend to CUDA, but the resolved protocol and frozen environment record must capture that override before any job is executed; backend changes during a run are invalid.

The default quota is $\alpha=0.95$, with a registered grid $\{0.85,0.90,0.95,0.99\}$. The $\alpha$ stage covers all seven datasets at seed 42. Ablations cover Madelon, ISOLET, Credit-Card-Fraud, and Abalone at seed 42. All other formal comparisons use the complete five-seed scope. Because the repository currently contains no completed formal campaign, the remainder of this section states questions, outcomes, and acceptance rules without result values.

\begin{table}[t]
\caption{Dataset summary. MI = mutual-information filtering ($p\leq100$ keep all, $p>100$ keep top-100).}
\label{tab:datasets}
\begin{center}
\footnotesize
\begin{tabular}{l r r r l}
\toprule
\textbf{Dataset} & \textbf{Native dim.} & \textbf{After MI} & \textbf{Samples} & \textbf{Task} \\
\midrule
Spambase & 57 & 56 & 4,601 & binary \\
Madelon  & 500 & 100 & 2,600 & binary \\
ISOLET   & 617 & 100 & 7,797 & 26-class \\
Gisette  & 5,000 & 100 & 7,000 & binary \\
Credit-Card & 29 & 29 & 284,807 & binary \\
Crimedata & 128 & 100 & 1,994 & regression \\
Abalone   & 8 & 8 & 4,177 & regression \\
\bottomrule
\end{tabular}
\end{center}
\end{table}

\subsection{Harmful-Feature Removal}
\label{sec:removal}

\textbf{RQ1: Does MinPPI removal retain predictive utility at a useful compression level?} For each dataset and seed, the selection plan estimates MinPPI at $\alpha=0.95$ and records the raw harmful set. Before the held-out test is accessed, that set must pass a validation non-inferiority gate: the retained subset's validation utility may be at most 0.005 below the full model for classification and at most 0.01 below it for regression. If the raw set fails, the frozen plan reduces the removal set according to the registered validation-gate procedure; the resulting accepted set, not the raw set, determines the matched deletion budget.

The removal stage compares the full model, MinPPI removal, SHAP-bottom removal, MinShap-bottom removal, Banzhaf-bottom removal, and random matched removal. Each non-full control removes exactly the accepted number of features for that dataset and seed; random matched removal has 30 frozen repeats. Powershap and ShapG are related methods but are not registered removal controls, so this study will not present them as entries in a formal ranking table. Primary removal outcomes are compression, selected and full test utility, and $\Delta U=U_{\text{selected}}-U_{\text{full}}$. Model-fit, coalition, cache, and runtime accounting are diagnostic outcomes rather than evidence of predictive superiority.

\subsection{$\alpha$-Sensitivity}
\label{sec:alpha}

\textbf{RQ2: How does the registered quota change compression and utility preservation?} At seed 42, the alpha stage evaluates $\alpha\in\{0.85,0.90,0.95,0.99\}$ on all seven datasets. It records harmful count, compression, validation utility, test utility, full-test utility, $\Delta U$, and the ratio of selected to full-test utility. These trajectories are descriptive operating points: the protocol does not assume monotonic harmful counts, monotonic compression, or universal preservation. The default $\alpha=0.95$ was frozen before the formal campaign and will not be chosen retrospectively from the sweep.

\subsection{Pre-Filtered Multi-Objective Search}
\label{sec:nsga2}

\textbf{RQ3: Does MinPPI prefiltering change NSGA-II's Pareto search under a matched evaluation budget?} The registered comparison is only NSGA-II-full versus NSGA-II with the validation-gated MinPPI prefilter. Both variants use population 40, 800 evaluations, and hypervolume reference point $(1.1,1.1)$; their search settings and evaluation protocol are otherwise identical. The recorded endpoints are validation hypervolume, held-out test hypervolume of the validation-selected front, convergence evaluation count, dimensional reduction, and complete fit/cache/runtime accounting. This protocol does not support a formal six-baseline rank table or critical-difference analysis.

\subsection{Estimator and Utility Ablations}
\label{sec:ablation}

\textbf{RQ4: Which estimator, guidance, and utility choices change the recovered MinPPI labels?} The estimator stage covers all seven datasets and all five canonical seeds, comparing exact, random MC, MI-correlation-guided, and hybrid estimates at the dataset-specific budget ladders; exact comparisons are interpreted only where exact enumeration is allowed. Separately, the registered seed-42 ablation scope covers Madelon, ISOLET, Credit-Card-Fraud, and Abalone. Its guidance axis compares random, mutual-information-only, correlation-only, and MI-correlation seeds, while its dispatcher axis compares exact, random MC, MI-correlation-guided, and hybrid execution. Their outcomes include harmful fraction and the estimator agreement metrics emitted by the canonical summary. On ISOLET alone, the utility axis compares accuracy with formal OVR AUC while keeping the probability-producing classifier fixed. Accuracy results are reported as ablation evidence and cannot replace the OVR-AUC formal endpoint.

\subsection{Pre-Registered Reporting and Acceptance Gates}
\label{sec:protocol}

Formal removal and NSGA-II summaries use the five seeds as the independent units. For each dataset, method, parameter, and metric, the canonical summary reports $n=5$, mean, sample standard deviation, median, minimum, and maximum. The 30 random-removal repeats are averaged within seed before the five-seed aggregation, so they do not inflate the inferential sample size. NSGA-II-full and MinPPI-prefilter are compared by the exact paired sign-flip test on aligned seeds, with paired rank-biserial effect size. Holm adjustment is applied separately within each registered endpoint family: validation hypervolume, test hypervolume, and convergence evaluation count. These low-powered analyses will be labeled descriptive; no result will be promoted to a broad method ranking.

Three gates precede any empirical claim. First, the frozen resolved protocol, source fingerprint, environment, manifest, and dataset identities must validate, and all reported rows must come from the formal rather than smoke run. Second, MinPPI removal must pass the pre-test validation non-inferiority gate described in Section~\ref{sec:removal}; held-out test values are reported outcomes, not inputs to subset choice. Third, the campaign summary must verify the complete canonical five-seed scope and every required upstream artifact and must mark the new provenance bundle as publication-eligible. If an approved omission exists, publication remains suppressed. Until all three gates pass, tables, figures, significance statements, runtime comparisons, and safety or superiority conclusions remain absent from the manuscript.

\section{Discussion and Conclusion}
\label{sec:discussion}

We introduced MinPPI, a game-theoretic framework that replaces the mean marginal contribution with the minimum under a binary utility-threshold voting game, yielding a principled three-way classification of features. The min--binary pairing occupies an empty cell in the design space of game-theoretic feature selectors and produces a discrete $\{-1,0,+1\}$ output requiring no arbitrary cutoff---unlike MinShap, whose continuous scores need a user-chosen threshold to decide which features to remove.

MinPPI's core contribution is the ensemble of three capabilities: (i)~threshold-free three-way classification via the min--binary pairing; (ii)~MI-correlation-guided adversarial search with a conditional geometric detection guarantee; and (iii)~$\alpha$-parameterized multi-objective selection via parametric scalarization. Unlike SHAP-based methods, which rank by mean marginal contribution and require a user-chosen $k$, MinPPI proposes a discrete candidate set that can be passed through a validation gate. Unlike MinShap, MinPPI's binary $v_\alpha$ eliminates an additional score cutoff. Unlike NSGA-II alone, MinPPI can provide structured prior information about harmful witnesses for evolutionary search.

The empirical contribution in this draft is the pre-registered protocol in Section~\ref{sec:experiments}. It asks whether validation-gated MinPPI removal retains task utility, how quota and estimator choices change the selected set, and whether prefiltering alters NSGA-II hypervolume or convergence under a matched budget. The present evidence state does not answer those questions: the formal results directory is empty, and no publication-eligible provenance bundle is available. Consequently, this version makes no claim of empirical safety, superiority, monotonic quota behavior, computational acceleration, or downstream search improvement. Those conclusions, if supported, must be written from a newly validated formal bundle using the registered five-seed summaries and acceptance gates.

\textbf{Limitations.} Theorem~2 is conditional on a harmful witness having positive probability under Algorithm~B's full perturbation distribution, an assumption that cannot be verified without ground-truth witnesses; it supplies no universal sample-complexity advantage over random sampling. On real data Algorithm~B remains a heuristic whose behavior must be measured by the registered ablations. The min estimator is conservatively biased upward when only a subset of coalitions is sampled (Lemma~1), so finite-budget estimates can miss harmful witnesses. The protocol uses XGBoost as the base learner and does not establish transfer to other model families. Finally, five seeds provide limited inferential power, which is why the registered statistical comparisons are explicitly descriptive.

\textbf{Future work.} The UTVG framework extends naturally to bipartite (feature $\times$ architecture) factored games, where cross-Shapley interaction values~\citep{grabisch1999axiomatic,rabitti2019shapley} could jointly select features and neural architectures. The min--binary pairing may also apply to adversarial robustness, where worst-case feature interactions under perturbation are the quantity of interest.

% ---- Appendix ----
\bibliography{icomp2026_conference}
\bibliographystyle{icomp2026_conference}

\appendix

\section{Proofs}
\label{app:proofs}

\begin{proof}[Proof of Property~1]
For any $S \subseteq N$, the marginal $\Delta_i(S) = v_\alpha(S \cup \{i\}) - v_\alpha(S)$ takes values in $\{-1, 0, +1\}$ because $v_\alpha \in \{0,1\}$. Both $\text{MinSS}_i$ and $\text{MinBz}_i$ are minima of subsets of $\{-1, 0, +1\}$, hence remain in $\{-1, 0, +1\}$.
\end{proof}

\begin{proof}[Proof of Property~2]
(a) $(\Rightarrow)$ $\text{MinSS}_i = +1$ means every permutation yields marginal $+1$, so $\forall S \subseteq N\setminus\{i\}: v_\alpha(S)=0, v_\alpha(S\cup\{i\})=1$. $(\Leftarrow)$ If the condition holds, every marginal is $+1$, so $\min = +1$. (b) $\text{MinSS}_i = -1$ iff $\exists \pi$ with $-1$ marginal, giving witness $S = S_\pi^i$. Conversely, any such $S$ yields $-1$ for every $\pi$ where $S$ is the prefix of $i$. (c) By Property~1, if not $+1$ or $-1$, the value must be $0$.
\end{proof}

\begin{proof}[Proof of Property~3]
$\text{MinSS}_i = \min_{\pi} \Delta_i(S_\pi^i)$ and $\text{MinBz}_i = \min_{S} \Delta_i(S)$. For every $S\subseteq N\setminus\{i\}$ there are $|S|!\,(n-|S|-1)!\geq1$ permutations in which $S$ is exactly the predecessor set of $i$. Hence the set of distinct predecessor coalitions is $2^{N\setminus\{i\}}$; multiplicities vary with $|S|$ but do not change a minimum. Both definitions therefore minimize over the same distinct marginal values.
\end{proof}

\begin{proof}[Proof of Lemma~1]
The estimate $\widehat{\text{MinPPI}}_i^{\text{MC}} = \min_{S \in \mathcal{S}_K} \Delta_i(S)$ satisfies $\min \mathcal{S}_K \geq \min \mathcal{M}_i$ deterministically because $\mathcal{S}_K \subseteq \mathcal{M}_i$. Hence the estimate is conservatively biased upward: it may miss a harmful signal but never falsely flag a neutral feature as harmful.
\end{proof}

\begin{proof}[Proof of Theorem~2]
For one complete perturbation draw $S\sim P_i$, the support assumption gives $\Pr[S=S_i^*]=\rho_i>0$. Because $S_i^*$ is harmful, the event $S=S_i^*$ implies detection. Thus one draw misses that particular witness with probability $1-\rho_i$. Independence of the $K_{\text{adv}}$ draws gives probability $(1-\rho_i)^{K_{\text{adv}}}$ of never sampling it, so $\Pr[\text{detection}]\geq1-(1-\rho_i)^{K_{\text{adv}}}$. The bound makes no claim when $\rho_i=0$ and no deterministic one-draw claim when the implementation perturbs rather than explicitly evaluates the seed coalition.
\end{proof}

\section{Additional Experimental Details}
\label{app:details}

\textbf{Protocol.} For each canonical seed, the split fractions are train 0.64, validation 0.16, and test 0.20, with stratification for classification. Variance filtering, mutual-information selection to at most 100 features, and standardization are fitted without test leakage. The base learner is XGBoost with a 100-round configured reference/cap, maximum depth 4, learning rate 0.1, and \texttt{n\_jobs=1}. Training/validation-only calibration selects and freezes a dataset-specific round count from $\{10,20,40,100\}$ before formal test evaluation; classification evaluation consumes predicted probabilities. Binary utilities are ROC AUC, ISOLET's formal utility is OVR AUC, and regression utilities are $R^2$. The five formal seeds are 42, 123, 456, 789, and 1024. MinPPI uses $\alpha=0.95$ by default and $\{0.85,0.90,0.95,0.99\}$ for the registered alpha stage.

The estimator methods are exact, random MC, MI-correlation-guided, and hybrid. The per-dataset budgets are $[2,5,10,25]$ except ISOLET, which uses $[2,5,10]$; exact enumeration is allowed only at 16 or fewer selected features. Removal controls are full, MinPPI, SHAP-bottom, MinShap-bottom, Banzhaf-bottom, and random matched, with 30 random repeats. Removal non-inferiority margins are 0.005 for classification and 0.01 for regression. NSGA-II compares only full and MinPPI-prefilter variants, with population 40, 800 evaluations, and reference point $(1.1,1.1)$. The alpha stage uses all datasets at seed 42; the ablation stage uses Madelon, ISOLET, Credit-Card-Fraud, and Abalone at seed 42, with the utility axis restricted to ISOLET.

The portable protocol defaults to CPU. A formal campaign may explicitly select CUDA only before run creation; the resolved protocol, model policy, frozen environment, and source fingerprint must record that selection. Resume rejects protocol drift. Smoke artifacts, diagnostic benchmarks, runtime forecasts, and screening outputs are not formal evidence and cannot populate the paper. A formal result becomes eligible for reporting only after the summary verifies the canonical five-seed scope, all required upstream artifacts, and a publication-eligible provenance bundle.

\textbf{Method design space.} Table~\ref{tab:op-table} classifies prior work and MinPPI along two axes: aggregation operator (mean vs.\ min) and value-function type (binary vs.\ continuous). In this comparison, MinPPI occupies the min $\times$ binary cell. Its discrete $\{-1,0,+1\}$ output follows from minimizing binary-game marginals; continuous-value methods instead retain the scale of their underlying value function.

\begin{table}[H]
\caption{Design space of game-theoretic feature importance methods. MinPPI occupies the empty cell---min operator $\times$ binary value function---producing a discrete three-way output.}
\label{tab:op-table}
\begin{center}
\begin{tabular}{lccc}
\toprule
\textbf{Method} & \textbf{Aggregation} & $v(S)$~\textbf{type} & \textbf{Output range} \\
\midrule
Shapley-Shubik~\citep{shapley1954method} & mean & binary $\{0,1\}$ & $\subseteq [-1,1]$ \\
Banzhaf~\citep{banzhaf1965weighted}       & mean & binary $\{0,1\}$ & $\subseteq [-1,1]$ \\
Shapley/SHAP~\citep{lundberg2017unified} & mean & continuous & \shortstack{real-valued\\scale-dependent} \\
MinShap~\citep{zheng2026minshap}         & min  & continuous & \shortstack{real-valued\\scale-dependent} \\
\midrule
\textbf{MinPPI (this work)} & \textbf{min} & \textbf{binary $\{0,1\}$} & $\mathbf{\{-1,\,0,\,+1\}}$ \\
\bottomrule
\end{tabular}
\end{center}
\end{table}

\textbf{ISOLET utility choice.} On ISOLET (26-class), the formal voting-game and evaluation utility is OVR AUC computed from class probabilities. Accuracy is evaluated only as the registered seed-42 utility ablation. The ablation will report any agreement or disagreement between the two label sets; the protocol does not assume equivalence.

\section{Estimator Verification Plan}
\label{app:verify}

Constructed voting games may be used as software checks for Properties~1--3, Lemma~1, and the conditional bound in Theorem~2. Such checks are not part of the formal seven-dataset evidence and will not be mixed with campaign tables or figures. Any numerical verification included in a submission must be regenerated from an identified artifact and linked to its provenance; this draft therefore states only the analytic proofs above and the pre-registered real-data estimator comparisons in Section~\ref{sec:ablation}.

\end{document}
